\chapter{Formal Structure}
\label{chap:formal}
%
In this chapter you will find basic information on the formal structure of your 
thesis.

\section{Sequence}
\label{sec:aufbau}
A scientific paper usually consists of the following parts:
%
\begin{enumerate}
 \item Cover sheet
 \item Abstract (optional)
 \item Table of contents
 \item List of figures and tables (also common at the end)
 \item List of abbreviations (also common at the end)
 \item Introduction
 \item Main part
 \item Summary\slash conclusions
 \item Bibliography
 \item Appendices (optional)
 \item Declaration
\end{enumerate}

\section{Cover Page}
The cover page includes: title of the thesis, type of thesis, author, 
matriculation number, submission date, supervisor and second reviewer. For 
theses longer than 15~pages, the cover page is included in the number of pages 
but not numbered.
%
\section{Table of Contents}
\label{sec:listOfContents}
We recommend a decimal structure as laid out in this document. If subheadings 
are used within a chapter, there must be at least two: where there is a 2.1, 
there must be a 2.2.

The table of contents always contains the page references for the listed items; 
however, it is not listed in the table of contents itself. The pages that the 
table of contents itself takes up can be counted in Roman numerals.

For a thesis, a structure of at least three levels is usually used. As a rule, 
only up to four levels are shown in the table of contents. However, you must 
find out about and consider the conventions in your subject.

\section{List of Figures and Tables}
Figures and tables are listed in corresponding directories. In this template, 
they appear directly after the table of contents. The corresponding pages can 
then be numbered in Roman numerals. However, the lists can also be placed at 
the end of the work before or after the bibliography. In this case, they are 
given regular page numbers. 

\section{List of Abbreviations}
The number of abbreviations should be kept to a minimum. The list of 
abbreviations only contains important subject-specific abbreviations in 
alphabetical order, in particular abbreviations of organisations, associations 
or laws. Common abbreviations such as i.\,a., e.\,g., etc. are not included.

For technical implementation with \LaTeX, see also \cref{sec:template}.

\section{Bibliography}
The bibliography is organized alphabetically by author name. It contains all 
sources cited in the text - and only these. Several writings by the same person 
are organized by year of publication. You must distinguish between writings by 
the same person from the same year of publication. In engineering, a number or 
author abbreviation is often placed before the name in square brackets. You can 
learn more about this and other important rules of citation in the Writing 
Centre's e-learning courses%
\footnote{\href{https://ilu.th-koeln.de/goto.php?target=cat\_52109\&client\_id=thkilu}{https://ilu.th-koeln.de/goto.php?target=cat\_52109\&client\_id=thkilu}}%
.

Appropriate software tools such as Citavi or Zotero, which are compatible with 
various word processing programs, are suitable for managing the literature 
used.

\section{Spelling, Grammar}
When submitting your work, make sure that your German or English is flawless. 
If errors impair readability, this can have a negative impact on your grade. It 
is therefore essential that you use the spell checker in your word processing 
programme, even if it does not detect all errors.

For those who feel unsure about this topic, we recommend the Writing Centre's 
e-learning courses%
\footnote{\href{https://ilu.th-koeln.de/goto.php?target=cat\_52109\&client\_id=thkilu}{https://ilu.th-koeln.de/goto.php?target=cat\_52109\&client\_id=thkilu}}%
. If necessary, please also contact the representative for students with 
disabilities%
\footnote{\href{https://www.th-koeln.de/studium/studieren-mit-beeintraechtigung\_169.php}{https://www.th-koeln.de/studium/studieren-mit-beeintraechtigung\_169.php}}
.

\section{Scope of the Thesis}
All subjects specify binding upper and lower limits that must generally be 
adhered to. Indexes and appendices are generally not counted. In individual 
cases~--~especially in the case of empirical work~--~deviating agreements can 
be made with the supervisor.
